\SourceLeaf{302}
\par\bigskip\pdfbookmark[0]{Remarks on On the Connection}{nexus-bemerkungen}
\begin{center}
\textsc{Remarks on the Paper}\\[1em]
\textsc{On the Connection between the Number of Classes into Which Binary Quadratic Forms Are Distributed and Their Determinant.}
\end{center}
\BookRule{.16\textwidth}
\par\medskip\pdfbookmark[1]{On I and II.}{nexus-note-I-II}\centerline{On I and II.}\smallskip

The second formula for the number of points inside the circle (I, Article 3, and II, Article 5) follows by considering the square inscribed in the circle, whose sides are parallel to the coordinate axes; comparing the two formulas leads to the theorem
\[
r'+r''+\ldots+r^{(q)}=qq+r^{(q+1)}+r^{(q+2)}+\ldots+r^{(r)}
\]
which is also easy to prove arithmetically. This in turn gives the correctness of the first of the following two rules, which were found on a separate sheet:

“Solutions of the equation $xx+yy\leqq A$; formula
\[
1+4\surd A+4\surd\tfrac12 A+8\Sigma(\surd(A-nn)-n)
\]
where the fractional part of each square root is dropped, and the sum is taken from $n=1$ to $n=\surd A$ ‘(should read $\surd\tfrac12 A$)’.

Another formula
\[
1+4\{A-\tfrac A3+\tfrac A5-\tfrac A7+\tfrac A9-\tfrac A{11}\ldots\}
\]
where the fractional part of each term is dropped.”

The latter formula follows from the theorem applied later (I, Article 6) concerning the number of all distinct representations of a given number by the form $xx+yy$ (compare Disqq. Arithm., Article 182, note). This theorem can easily be put in the following form: the number of distinct representations of a positive integer $m$ by the form $xx+yy$ is $=4(a-b)$, where $a,b$ denote the numbers of divisors of $m$ of the respective forms $4n+1,\ 4n+3$. By comparing
\SourceLeaf{303}
this arithmetical formula with the average number of representations obtained by geometrical considerations (in I, Article 5, or II, Article 4), one easily obtains, with full rigor, the familiar result
\[
\frac\pi4=1-\tfrac13+\tfrac15-\tfrac17+\ldots
\]
which is derived in the paper (I, Article 7) by a similar comparison, but with the help of infinite products.
\BookRule{.25\textwidth}
